\documentclass[a4paper,10pt]{article}
\usepackage[utf8]{inputenc}
\usepackage[T1]{fontenc}
\usepackage{lmodern}
\usepackage[ngerman]{babel}
\usepackage[a4paper,portrait,left=22mm,right=22mm,top=22mm,bottom=24mm]{geometry}
\usepackage{longtable,array,booktabs}
\usepackage{xcolor}
\usepackage{enumitem}
\usepackage{microtype}
\usepackage{fancyhdr}
\usepackage{amsmath}
\usepackage[font=small,skip=5pt]{caption}
\pagestyle{fancy}
\fancyhf{}
\fancyfoot[L]{\small\texttt{C3\_SPM\_Blatt3.pdf}}
\fancyfoot[R]{\small Seite \thepage}
\renewcommand{\headrulewidth}{0pt}
\renewcommand{\footrulewidth}{0.4pt}
\setlength{\footskip}{22pt}
\setlength{\parindent}{0pt}
\setlength{\parskip}{4pt}
\setlength{\emergencystretch}{3em}
\setlist{nosep,leftmargin=*,topsep=3pt}
\clubpenalty=10000
\widowpenalty=10000
\displaywidowpenalty=10000
\raggedbottom
\setlength{\tabcolsep}{4pt}
\renewcommand{\arraystretch}{1.15}
\setlength{\LTpre}{4pt}
\setlength{\LTpost}{8pt}
\setlength{\LTcapwidth}{\textwidth}
\newcommand{\sectionspace}[1]{\par\ifdim
                                      \dimexpr\pagegoal-\pagetotal\relax<#1\relax\newpage
\fi\addvspace{14pt}}
\newcommand{\offen}{\emph{noch offen}}


\title{Projektziele, Projektstruktur-, Ablauf-, Termin- und Kostenplan\\{\large Übungsblatt 3}}
\author{Gruppe C3 / MoinMetrics\\{}Alexander Borrmann\\{}Lukas Pancratius\\{}Ino Wölke}
\date{Hochschule Emden-Leer\\Software-Projektmanagement\\Wintersemester 2026/27\\Prof. Dr. Gerrit Jan Veltink\\[1em]Abgabedatum: noch offen}
\begin{document}
\maketitle
\thispagestyle{fancy}
\clearpage
\section*{Aufgabe 1: Projektziele als Voraussetzung für die Planung}
% Ober- und Unterziele formulieren, Muss- und Soll-Ziele kennzeichnen.
% Auf die Bezeichner aus dem Anforderungsdokument verweisen.
\offen

\section*{Aufgabe 2: Projektstrukturplan}
% Die vom Projektteam zu erledigenden Arbeitspakete hierarchisch darstellen.
\offen

\section*{Aufgabe 3: Projektablaufplan}
\subsection*{a) Vorgangsknotennetz ohne Zeitschätzungen}
% Alle Abhängigkeiten der Arbeitspakete berücksichtigen.
\offen

\subsection*{b) Ausführlicher Netzplan}
% Dauer realistisch schätzen; vier Zeitangaben und vier Pufferzeiten berechnen.
% Den kritischen Pfad markieren.
\offen

\subsection*{c) Excel-Datei zum Netzplan}
% netzplan.xlsx als Vorlage nutzen; sämtliche Spalten mit Werten, ohne Formeln.
% Datei netzplan_C3.xlsx nennen und mit netzplan_check.py prüfen.
\offen

\section*{Aufgabe 4: Terminplan}
% Anfangs- und Endzeitpunkte datieren; Wochenenden, Feiertage und Ferien beachten.
\offen

\section*{Aufgabe 5: Kostenplan}
% Projektstunden aus der Aufwandsschätzung mit realistischem Stundensatz bewerten.
\offen
\end{document}
